\chapter{Sécurité}\label{ch:pl:security}

Le 21 février 2025, une grande plateforme d'échange de cryptoactifs a signé une transaction issue de l'un de ses portefeuilles, qui exigeait plusieurs signatures. Les signataires
ont approuvé ce que leurs écrans leur montraient~: une transaction légitime. L'application web qui l'avait préparée avait été modifiée deux jours plus tôt pour
afficher des données de transaction légitimes tout en changeant ce qui était réellement signé, et les signatures ont donné aux attaquants l'éther du portefeuille. Le
FBI a attribué le vol d'environ 1.5~milliard de dollars à la Corée du Nord. Aucune cryptographie n'a été cassée~: les attaquants ont changé ce qu'on montrait aux personnes.
La sécurité d'une société de trading consiste surtout à déterminer qui et quoi a le droit d'agir, et à s'assurer que ce qu'on voit est bien ce qu'on signe. Ce chapitre
écrit la politique d'accès de la société sous forme de données et la vérifie, garde les secrets et les clés des plateformes sous des règles strictes, segmente le réseau,
et traque l'initié qui copie un peu trop chaque jour.

\section{Ce qu'une société de trading protège, et contre qui}

\begin{definition}[Modèle de menace]\label{def:pl:security:threat}
Un \emph{modèle de menace}\index{modèle de menace} recense ce qu'un système doit protéger (ses actifs), contre qui (les adversaires et leurs moyens), par quels chemins (la
surface d'attaque) et quels contrôles couvrent chaque chemin, afin que l'effort de sécurité porte là où les pertes seraient les plus lourdes.
\end{definition}

Les actifs d'une société de trading sont inhabituels à deux égards. Certains sont de l'argent qui se déplace de lui-même sur une signature --- soldes sur les plateformes, portefeuilles (livre~3,
chapitre~24), instructions de paiement --- de sorte qu'un identifiant volé représente une perte en quelques minutes. D'autres sont de l'information dont la valeur tient à son secret ---
code de stratégie, signaux, paramètres, positions --- de sorte qu'une copie constitue une perte même si rien n'est supprimé. Les adversaires se répartissent ainsi~: criminels et
États pour les premiers, concurrents et employés partants pour les seconds. Le risque opérationnel du livre~6, chapitre~28, couvre les deux~; ce
chapitre construit les contrôles.

\section{Identité, accès et moindre privilège}

\begin{definition}[Moindre privilège, contrôle d'accès par rôle]\label{def:pl:security:rbac}
Le \emph{moindre privilège}\index{moindre privilège} n'accorde à chaque personne et à chaque programme que les permissions nécessaires à sa tâche, et seulement pendant
le temps où elles lui sont nécessaires. Le \emph{contrôle d'accès par rôle}\index{contrôle d'accès par rôle} (RBAC) accorde des permissions --- une action sur une ressource --- à des rôles, et des rôles
à des personnes, de sorte que l'accès suit la fonction plutôt que l'individu.
\end{definition}

\begin{definition}[Authentification multifacteur]\label{def:pl:security:mfa}
L'\emph{authentification multifacteur}\index{authentification multifacteur} exige plus d'un type distinct de facteur d'authentification --- quelque chose que l'on connaît, que l'on possède
ou que l'on est --- avant l'accès~; les recommandations actuelles du NIST exigent, à leur niveau d'assurance intermédiaire, deux facteurs distincts et la
mise à disposition d'un facteur résistant au hameçonnage.
\end{definition}

La société du chapitre compte douze rôles et cent personnes. Sa politique est une donnée (\texttt{firm.accessctl})~: les rôles accordent des couples (ressource, action), les personnes
détiennent des rôles, les règles de refus l'emportent sur les autorisations, et cinq contraintes de séparation des tâches (livre~6, chapitre~28) nomment des couples de
permissions qu'aucune personne seule ne doit détenir.

\omcode{code/platforms/29-security/python/pl_accessctl.py}{22}{27}{Les cinq contraintes de séparation des tâches de la société, chacune étant un couple de
permissions qu'une seule personne ne doit pas détenir.}\label{lst:pl:security:sod}

\omcode{code/firm/accessctl/firm_accessctl.py}{34}{56}{Le modèle de politique~: permissions obtenues par les rôles, moins les refus~; combinaisons toxiques repérées au niveau
des rôles et des personnes.}\label{lst:pl:security:policy}

\begin{table}[htbp]
\centering\small
\begin{tabular}{lll}
\toprule
rôle & première permission & seconde permission\\
\midrule
développeur & fusionner le code & déployer en production\\
responsable de desk & proposer les paramètres & approuver les paramètres\\
opérations & créer les paiements & libérer les paiements\\
fiabilité des sites & créer les clés de plateforme & retirer des fonds des portefeuilles\\
trader & soumettre des ordres & modifier les transactions\\
\bottomrule
\end{tabular}
\caption{Les cinq rôles qui, avant correction, détenaient à eux seuls les deux moitiés d'une contrainte de séparation des tâches. Par leur intermédiaire et par celui de
trois chercheurs qui étaient aussi gestionnaires des risques, 60 des 100 personnes détenaient une combinaison toxique~; après la scission des rôles, l'ajout d'un responsable
des mises en production et la suppression des seconds rôles, aucune ne la détient.}\label{tab:pl:security:toxic}
\end{table}

Chaque rôle toxique a été créé par commodité, et chacun correspond à un type de perte connu~: un trader qui modifie ses propres saisies peut dissimuler des pertes
(les cas de fraude de trader du livre~6)~; un développeur qui déploie sa propre fusion contourne la revue du chapitre~\ref{ch:pl:testing-and-continuous-delivery}~;
un responsable de desk qui approuve ses propres paramètres contourne le principe des quatre yeux du chapitre~\ref{ch:pl:signal-serving-and-parameter-management}. Le contrôle tient en
quelques lignes de code appliquées aux données de politique et doit figurer dans l'\omterm{def:pl:testing-and-continuous-delivery:ci}{intégration continue} de la politique~: une modification qui crée une combinaison toxique échoue
comme un test en échec.

\section{Secrets et clés}

\begin{definition}[Gestion des secrets, rotation des secrets]\label{def:pl:security:secrets}
La \emph{gestion des secrets}\index{gestion des secrets} conserve les identifiants --- mots de passe, clés d'API, clés privées, mots de passe de bases de données --- dans un
coffre qui les délivre aux programmes authentifiés sous forme de baux et qui journalise chaque accès, et jamais dans le code ni dans les fichiers de configuration. La
\emph{rotation des secrets}\index{rotation des secrets} remplace un secret selon un calendrier et à tout soupçon, de sorte qu'un secret divulgué n'est valide que jusqu'à sa
prochaine rotation~: c'est sa période d'utilisation, au sens du NIST.
\end{definition}

\begin{definition}[Module matériel de sécurité]\label{def:pl:security:hsm}
Un \emph{module matériel de sécurité}\index{module matériel de sécurité} (HSM) est un dispositif résistant à la falsification qui génère et conserve des clés et
effectue des opérations avec elles, de sorte que les clés privées ne sortent jamais de lui sous forme claire~; la signature s'effectue à l'intérieur, sur les données que l'hôte lui présente.
\end{definition}

La rotation borne l'exposition~: un secret qui fuit à un moment quelconque et qui est renouvelé tous les 90~jours reste valide 45~jours en moyenne~; renouvelé chaque semaine,
trois jours et demi. Un HSM, ou le calcul multipartite du livre~3, chapitre~24, retire la clé de l'hôte --- mais non la question que pose
le hameçonnage~: le dispositif signe ce qu'on lui montre. La parade consiste à rendre ce qui est signé vérifiable indépendamment de l'écran qui l'a préparé~:
un second canal qui affiche le contenu réel de la transaction, un moteur de politique qui refuse un type de transaction (une mise à niveau de contrat, une nouvelle
destination) quelle que soit l'approbation des signataires.

\begin{omfigure}
\begin{tikzpicture}[font=\small,node distance=6mm,
  box/.style={draw=omDef,fill=omDef!8,rounded corners=2pt,align=center,minimum height=11mm,text width=31mm},
  bad/.style={draw=omThm,fill=omThm!10,rounded corners=2pt,align=center,minimum height=11mm,text width=24mm},
  arr/.style={-{Stealth[length=2mm]},thick}]
\node[bad] (app) {application web\\(modifiée)};
\node[box,right=26mm of app,yshift=10mm] (scr) {écran du signataire~:\\le virement attendu};
\node[box,right=26mm of app,yshift=-10mm] (dev) {dispositif de signature~:\\une autre transaction};
\node[box,right=10mm of dev,text width=27mm] (chk) {décodeur indépendant\\et moteur de politique};
\node[box,right=16mm of chk,text width=14mm] (chain) {porte\-feuille};
\draw[arr] (app) -- node[above,sloped,font=\scriptsize]{ce qui est montré} (scr);
\draw[arr,omThm] (app) -- node[below,sloped,font=\scriptsize]{ce qui est signé} (dev);
\draw[arr] (dev) -- (chk);
\draw[arr] (chk) -- node[below,font=\tiny,align=center]{seulement si\\autorisé} (chain);
\draw[dashed,gray] (scr.east) -| (chk.north);
\node[gray,font=\scriptsize,anchor=south west] at (scr.north east) {le contrôle effectué sur l'écran};
\end{tikzpicture}

\omcaption{L'attaque du hameçonnage en une image~: l'application modifiée montrait aux signataires une transaction et envoyait à leurs dispositifs une autre. Un
contrôle qui décode ce qui doit réellement être signé, indépendamment de l'application, et une politique qui refuse ce qu'aucun signataire ne devrait approuver
(une modification du contrat propre au portefeuille, une nouvelle destination) contrôlent la chose signée plutôt que la chose montrée.}\label{fig:pl:security:signing}
\end{omfigure}

\section{Identifiants des bourses et des portefeuilles}

\begin{definition}[Liste blanche de retraits]\label{def:pl:security:allowlist}
Une \emph{liste blanche de retraits}\index{liste blanche de retraits} (withdrawal allow-list) restreint les destinations vers lesquelles un compte ou une clé peut envoyer des actifs à une liste
fixée à l'avance, modifiée uniquement par une procédure distincte et plus lente, de sorte qu'un identifiant volé ne puisse pas transférer d'actifs au voleur.
\end{definition}

Les clés d'API des plateformes (livre~3, chapitre~15) sont les secrets les plus exposés de la société~: elles résident sur les serveurs de trading et signent chaque ordre. Dans ce chapitre, les
clés portent des périmètres d'action --- une clé de trading peut consulter et négocier mais pas retirer~; une clé de retrait ne peut que retirer, vers des adresses figurant sur la liste blanche --- et
chaque requête est signée avec le signataire HMAC du livre~13 et horodatée. Les contrôles de la plateforme déterminent ce qu'une clé volée peut faire
(le \cref{tab:pl:security:keys}).

\omcode{code/firm/accessctl/firm_accessctl.py}{99}{114}{Les contrôles de la plateforme sur une requête~: une clé connue, une signature valide, un horodatage récent, une
action dans le périmètre de la clé, et une adresse de retrait figurant sur la liste blanche.}\label{lst:pl:security:verify}

\begin{table}[htbp]
\centering\small
\begin{tabular}{lll}
\toprule
tentative & acceptée & motif\\
\midrule
négociation avec la clé de trading & oui & ok\\
retrait avec la clé de trading & non & action hors du périmètre de la clé\\
retrait vers une adresse extérieure & non & adresse absente de la liste blanche\\
retrait vers le dépositaire & oui & ok\\
signature falsifiée & non & signature invalide\\
rejouée une minute plus tard & non & requête périmée\\
\bottomrule
\end{tabular}
\caption{Une clé de trading volée et une clé de retrait volée soumises aux contrôles de la plateforme. Le voleur peut négocier --- ce qui peut encore coûter de l'argent
--- mais ne peut transférer des actifs vers aucune destination que la société n'a pas inscrite à l'avance.}\label{tab:pl:security:keys}
\end{table}

\section{Segmentation et risque interne}

\begin{definition}[Segmentation du réseau, architecture zero trust]\label{def:pl:security:segmentation}
La \emph{segmentation du réseau}\index{segmentation du réseau} (network segmentation) découpe le réseau d'une société en zones --- trading, recherche, bureautique, l'extérieur ---
avec des flux entre zones autorisés uniquement sur des chemins déclarés, de sorte qu'une compromission dans une zone n'atteigne pas les autres. Une \emph{architecture
zero trust}\index{architecture zero trust} (zero-trust architecture) accorde aucune confiance implicite du seul fait d'être à l'intérieur d'un réseau~: chaque requête est authentifiée et
autorisée séparément, sur la base de l'identité et du dispositif, avant l'ouverture d'une session vers une ressource.
\end{definition}

\begin{omfigure}
\begin{tikzpicture}[font=\small,
  zone/.style={draw=omDef,fill=omDef!6,rounded corners=3pt,minimum width=30mm,text width=30mm,minimum height=16mm,align=center},
  hot/.style={draw=omThm,fill=omThm!8,rounded corners=3pt,minimum width=30mm,text width=30mm,minimum height=16mm,align=center},
  arr/.style={-{Stealth[length=2mm]},thick}]
\node[hot] (t) at (0,0) {zone trading\\\scriptsize stratégies, passerelles, clés};
\node[zone] (r) at (6.2,0) {zone recherche\\\scriptsize base de ticks, notebooks};
\node[zone] (c) at (12.2,0) {zone bureautique\\\scriptsize courriel, ordinateurs portables};
\node[zone] (v) at (0,-2.8) {plateformes\\\scriptsize interconnexions};
\node[zone] (i) at (12.2,-2.8) {internet};
\draw[arr] (t) -- node[right,font=\scriptsize]{ordres uniquement} (v);
\draw[arr] ([yshift=3mm]t.east) -- node[above,font=\scriptsize]{exécutions, positions} ([yshift=3mm]r.west);
\draw[arr,dashed] ([yshift=-3mm]r.west) -- node[below,font=\scriptsize,align=center]{paramètres, par\\revue (chap.~16)} ([yshift=-3mm]t.east);
\draw[arr] (c) -- node[above,font=\scriptsize,align=center]{accès\\authentifié} (r);
\draw[arr] (i) -- node[right,font=\scriptsize]{courriel, web} (c);
\draw[omThm,thick] (6.2,-2.8) node[cross out,draw,minimum size=4mm]{} node[below=3mm,font=\scriptsize,align=center]{aucun chemin depuis internet\\ou depuis les ordinateurs portables\\vers le trading};
\end{tikzpicture}

\omcaption{Les zones de réseau du chapitre et les seuls chemins autorisés entre elles. La zone de trading, qui détient les clés d'API, dialogue avec les plateformes
et envoie les exécutions à la recherche~; les paramètres ne reviennent que par le chemin soumis à revue du chapitre~16~; rien en provenance de la zone bureautique ou d'internet
ne l'atteint directement.}\label{fig:pl:security:zones}
\end{omfigure}

\begin{definition}[Menace interne, prévention des fuites de données]\label{def:pl:security:insider}
Une \emph{menace interne}\index{menace interne} (insider threat) est un préjudice causé par une personne ayant un accès légitime --- un salarié, un prestataire --- par le vol,
sabotage ou négligence. La \emph{prévention des fuites de données}\index{prévention des fuites de données} (data-loss prevention) est l'ensemble des contrôles qui détectent et bloquent la circulation
de données protégées hors des lieux où elles sont autorisées à se trouver~: copies sur supports amovibles, téléversements, courriels et volumes de lecture inhabituels.
\end{definition}

Le cas le plus connu du trading est aussi un avertissement sur les limites du droit~: en 2009, un programmeur a téléversé plus de 500\,000 lignes du code de trading à haute
fréquence de son employeur vers un serveur en Allemagne le dernier jour~; sa condamnation fédérale a été infirmée en appel en 2012, sur la portée des textes plutôt que sur les faits.
Une société ne peut pas compter sur les tribunaux (livre~16, chapitre~11, sur les secrets d'affaires)~; elle doit repérer
la copie. L'initié le plus difficile à repérer n'est pas celui qui prend tout le dernier jour, mais celui qui prend chaque jour un peu plus que d'habitude.

Le modèle du chapitre donne à chacune des cent personnes une année de journaux d'accès, soit 250 jours de trading --- les fichiers de code de stratégie que chacune lit
par jour, à son niveau propre, avec un rythme hebdomadaire et du bruit --- et introduit un initié qui, à partir d'un jour tiré au hasard, lit une part déclarée
de plus que d'habitude, chaque jour. Quatre détecteurs notent chaque personne par rapport à sa propre référence, la médiane et l'écart absolu médian sur 30 jours~:
un seuil journalier, ou une CUSUM des notes quotidiennes (livre~7, chapitre~13), chacun avec une référence qui se termine aujourd'hui ou une qui s'achève 60 jours
plus tôt. Les quatre sont calibrés sur des personnes non concernées, à raison d'une fausse alarme par mois pour cent personnes.

\omcode{code/firm/accessctl/firm_accessctl.py}{118}{137}{Notes robustes par personne par rapport à une référence qui peut être décalée, et la CUSUM unilatérale
qui cumule les petits excès quotidiens.}\label{lst:pl:security:scores}

\begin{omfigure}
\begin{tikzpicture}
\begin{axis}[width=11.5cm,height=5.6cm,xlabel={lecture supplémentaire par jour (\% au-dessus du niveau habituel de chacun, échelle logarithmique)},ylabel={initiés signalés (sur 20)},
  grid=major,grid style={gray!20},tick label style={font=\small},label style={font=\small},table/col sep=comma,
  xmode=log,xmin=4,xmax=120,ymin=0,ymax=21,xtick={5,10,20,30,50,100},xticklabels={5,10,20,30,50,100}]
\addplot[omThm,thick,mark=*] table[x=extra_pct,y=cusum_lagged_found]{figdata/platforms/29-security/detection.csv};
\addplot[omThm,thick,dashed,mark=o] table[x=extra_pct,y=cusum_trailing_found]{figdata/platforms/29-security/detection.csv};
\addplot[omDef,thick,mark=square*] table[x=extra_pct,y=day_lagged_found]{figdata/platforms/29-security/detection.csv};
\addplot[omDef,thick,dashed,mark=square] table[x=extra_pct,y=day_trailing_found]{figdata/platforms/29-security/detection.csv};
\node[omThm,font=\scriptsize,anchor=south east] at (axis cs:28,17) {CUSUM, référence décalée};
\node[omThm,font=\scriptsize,anchor=north west] at (axis cs:32,6) {CUSUM, référence glissante};
\node[omDef,font=\scriptsize,anchor=south east] at (axis cs:48,12.5) {un jour, référence décalée};
\node[omDef,font=\scriptsize,anchor=south] at (axis cs:45,1.6) {un jour, référence glissante};
\end{axis}
\end{tikzpicture}

\omcaption{Salariés introduits et signalés, sur vingt, selon la lecture quotidienne supplémentaire par rapport à leur habitude, à raison d'une fausse alarme par mois pour
cent personnes. Un seuil journalier sur une référence glissante ne les voit presque jamais~; une CUSUM sur une référence décalée de 60 jours en trouve 12 sur 20 à
+20~\% par jour et 19 sur 20 à +30~\%. Données~: \texttt{fig\_accessctl.py}.}\label{fig:pl:security:detection}
\end{omfigure}

Deux points comptent, et la \cref{fig:pl:security:detection} les sépare. Une seule journée où l'on lit 30~\% de plus que d'habitude se situe bien dans les fluctuations
ordinaires d'une personne~; seule l'accumulation --- la CUSUM --- transforme trente jours de ce type en signal. Et une référence glissante apprend la nouvelle habitude~: au bout de
trente jours, la copie de l'initié devient la normale de la personne, et le score revient à zéro. Décaler la référence de soixante jours empêche la copie de rester dans sa propre référence
assez longtemps pour être vue. Avec les deux, un initié qui lit 30~\% de plus que d'habitude est signalé en médiane au bout de 22 jours (19 sur
20), celui qui lit 50~\% de plus l'est en 12 (les 20)~; un seuil journalier sur une référence glissante n'en signale qu'un sur vingt.

\begin{omfigure}
\begin{tikzpicture}
\begin{axis}[width=11.5cm,height=5.2cm,xlabel={jour de négociation},ylabel={CUSUM des scores quotidiens},
  grid=major,grid style={gray!20},tick label style={font=\small},label style={font=\small},table/col sep=comma,
  xmin=90,xmax=250,ymin=0,ymax=60]
\addplot[omThm,thick] table[x=day,y=cusum_lagged]{figdata/platforms/29-security/insider.csv};
\addplot[omThm,dashed] table[x=day,y=cusum_trailing]{figdata/platforms/29-security/insider.csv};
\addplot[gray,thin] table[x=day,y=clean_cusum_lagged]{figdata/platforms/29-security/insider.csv};
\draw[black,dashed] (axis cs:90,12.709) -- (axis cs:250,12.709) node[pos=0.98,above left,font=\scriptsize]{seuil 12.7};
\draw[black,dotted] (axis cs:120,0) -- (axis cs:120,40) node[pos=0.75,left,font=\scriptsize,align=right]{début de la\\copie};
\node[omThm,font=\scriptsize,anchor=west] at (axis cs:205,57) {référence décalée};
\node[omThm,font=\scriptsize,anchor=west] at (axis cs:175,6) {référence glissante};
\end{axis}
\end{tikzpicture}

\omcaption{Un salarié malveillant introduit lit 30~\% de code de stratégie en plus que d'habitude à partir du jour 120. La CUSUM sur la référence décalée franchit son seuil
calibré au jour 131~; sur la référence glissante, elle n'atteint son propre seuil (11.2) qu'au jour 147, puis redescend une fois que la copie est entrée
dans la référence. La CUSUM d'une personne
saine (en gris) reste proche de zéro. Données~: \texttt{fig\_accessctl.py}.}\label{fig:pl:security:insider}
\end{omfigure}

\begin{dated}{2026-09}{Les affaires et les normes}\label{dat:pl:security:sources}
Le communiqué de presse du FBI du 26 février 2025 attribue à la Corée du Nord le vol d'environ 1.5~milliards de dollars d'actifs virtuels auprès de la
bourse Bybit vers le 21 février~; l'enquête publiée par Sygnia décrit un JavaScript modifié le 19 février dans l'application web Safe\{Wallet\}, de
sorte que son interface affichait des données de transaction légitimes alors que la transaction elle-même était altérée. Dans \emph{United States
v.\ Aleynikov} (2d Cir.\ 2012), la cour a annulé une condamnation pour le téléversement de plus de 500\,000 lignes de code source de trading. Le
document NIST SP 800-207 (2020) définit l'\omterm{def:pl:security:segmentation}{architecture zero trust}~; SP 800-63B (révision d'août 2025) les niveaux d'assurance de l'authentification~; SP 800-57 Part 1 (2020) la
durée de vie cryptographique d'une clé.
\end{dated}

\section{Tutoriel~: la copie lente}\label{tut:pl:security:tut}

\begin{tutorial}
\textbf{Objectif.} Vérifier la politique d'accès pour détecter les combinaisons toxiques, émettre et tester des clés de plateforme à portée limitée, et repérer un initié lent à un taux fixe de fausses
alarmes. \textbf{État final~:} \cref{tab:pl:security:toxic,fig:pl:security:detection}.
\begin{tutsteps}
\item \textbf{Politique}~: \texttt{pl\_accessctl.policies()} et \texttt{toxic\_summary()}, avant et après la correction.
\item \textbf{Clés}~: \texttt{stolen\_key\_attempts()}~; ajouter une clé avec un périmètre de négociation et une liste blanche, puis la tester.
\item \textbf{Coffre}~: \texttt{firm\_accessctl.Vault} avec un bail de 90 jours~; effectuer une rotation et lister ce qui arrive à échéance.
\item \textbf{Journaux}~: \texttt{access\_logs()} et \texttt{statistics(x)}~; \texttt{thresholds()}.
\item \textbf{Initié}~: \texttt{delays(extra)} pour chaque action de lecture supplémentaire.
\end{tutsteps}
\textbf{Pour aller plus loin.} Évaluer les personnes par rapport à leurs pairs du même rôle, en plus de leur propre historique~; ajouter des variables au niveau des
fichiers (jamais lus auparavant) et mesurer de combien l'initié est signalé plus tôt.
\end{tutorial}

\section{Construction~: contrôle d'accès}\label{bld:pl:security:accessctl}

\begin{build}
\bfield{Rôle} Accorder à chaque personne et à chaque programme uniquement ce dont son travail a besoin, garder les secrets et les clés de plateforme sous des règles qui bornent une fuite, et repérer une
copie lente à un taux de fausses alarmes annoncé.
\bfield{Interface} \texttt{Policy} (\texttt{allowed}, \texttt{permissions}, \texttt{toxic\_roles}, \texttt{toxic\_users}), \texttt{Vault},
\texttt{VenueKey}, \texttt{request}, \texttt{verify}, \texttt{robust\_scores}, \texttt{cusum}, \texttt{onsets}, \texttt{calibrate},
\texttt{first\_alarm}.
\bfield{Règles} Politique traitée comme des données relues, vérifiée contre les combinaisons toxiques en \omterm{def:pl:testing-and-continuous-delivery:ci}{intégration continue}~; secrets uniquement issus du coffre, sous bail,
soumis à rotation~; clés de plateforme à portée limitée, retraits uniquement vers des adresses figurant sur la liste blanche~; chaque requête signée et horodatée~; détecteurs calibrés sur des
populations saines et revus chaque mois.
\bfield{Tests d'acceptation} \texttt{code/firm/accessctl/tests/}~: autorisations, refus, rôles et utilisateurs toxiques~; baux du coffre, rotation et expiration~; chacune des vérifications
de la plateforme~; scores robustes décalés, CUSUM, calibration à un taux de fausses alarmes, et détection d'un décalage persistant.
\bfield{Pour aller plus loin} Accès just-in-time accordés puis expirés~; notation par groupes de pairs~; décodeur de transactions indépendant pour les signataires.
\end{build}

\begin{omsources}
\item Federal Bureau of Investigation, \emph{North Korea Responsible for \$1.5 Billion Bybit Hack}, PSA 250226, 26 février 2025~; Sygnia,
enquête sur le piratage de Bybit, 2025.
\item \emph{United States v.\ Aleynikov}, 676 F.3d 71 (2d Cir.\ 2012).
\item NIST SP 800-207, \emph{Zero Trust Architecture}, 2020~; SP 800-63B, \emph{Authentication and Authenticator Management}~; SP 800-57 Part 1
Rév.~5, \emph{Recommendation for Key Management}, 2020.
\item le livre~3, chapitres~14, 15 et~24 (clés, API de plateforme, portefeuilles)~; le livre~6, chapitre~28 (séparation des tâches)~; le livre~16, chapitre~11
(secrets d'affaires).
\end{omsources}

\section{Exercices}

\begin{exercise}[$\star$]\label{exo:pl:security:1}
Pourquoi la vérification des combinaisons toxiques regarde-t-elle les personnes autant que les rôles~?
\end{exercise}

\begin{exercise}[$\star$]\label{exo:pl:security:2}
Un secret fuit à un moment aléatoire. Combien de temps reste-t-il valide en moyenne avec une rotation de 90 jours, puis avec une rotation hebdomadaire~?
\end{exercise}

\begin{exercise}[$\star$]\label{exo:pl:security:3}
Que peut faire un voleur de la clé de négociation volée de l'entreprise, et que ne peut-il pas faire~?
\end{exercise}

\begin{exercise}[$\star\star$]\label{exo:pl:security:4}
Pourquoi une référence qui retarde sur le présent cesse-t-elle de voir une copie lente au bout d'environ un mois~?
\end{exercise}

\begin{exercise}[$\star\star$]\label{exo:pl:security:5}
Pourquoi un seuil journalier unique est-il le mauvais outil pour une copie lente, quel que soit son niveau~?
\end{exercise}

\begin{exercise}[$\star\star$]\label{exo:pl:security:6}
Les signataires du hook ont chacun vérifié la transaction sur leur écran. Quel contrôle aurait détecté l'attaque~?
\end{exercise}

\begin{exercise}[$\star\star\star$]\label{exo:pl:security:7}
\emph{Programmation.} Relancer \texttt{delays(0.3)} avec la valeur de référence \texttt{K} de la CUSUM fixée à 0.25 au lieu de 0.5 (en recalibrant les seuils).
Que se passe-t-il pour les détections, et pourquoi~?
\end{exercise}

\begin{exercise}[$\star\star\star$]\label{exo:pl:security:8}
\emph{Trouver la faille.} «~Notre réseau de trading est derrière le pare-feu, donc les systèmes qui s'y trouvent peuvent se faire confiance.~»
\end{exercise}

\section{Problème~: La copie lente}

\begin{problem}[{Problème du week-end --- la copie lente}]\label{pb:pl:security:1}
La politique, les clés, le coffre et les journaux d'accès du chapitre.

\medskip\noindent\textbf{Partie I --- Accès.}
\begin{enumerate}
\item Que protège la société, et contre qui~?
\item Quelles sont les cinq contraintes de séparation des tâches~?
\item Quels rôles étaient toxiques avant la correction, et combien de personnes détenaient une combinaison toxique~?
\item Qu'a changé la correction~?
\item Où la vérification des combinaisons toxiques doit-elle s'exécuter~?
\end{enumerate}

\medskip\noindent\textbf{Partie II --- Secrets et clés.}
\begin{enumerate}[resume]
\item Qu'apporte un coffre à un secret placé dans un fichier de configuration~?
\item Comment la rotation borne-t-elle une fuite~?
\item Quels sont les contrôles de la plateforme sur une requête~?
\item Qu'autorise chaque clé volée~?
\item Que protège un HSM, et que ne protège-t-il pas~?
\end{enumerate}

\medskip\noindent\textbf{Partie III --- Initiés.}
\begin{enumerate}[resume]
\item Comment construit-on et calibre-t-on les quatre détecteurs~?
\item Que signale chacun à +30~\% par jour~?
\item Pourquoi le CUSUM est-il important~?
\item Pourquoi le retard de la référence est-il important~?
\item Quand le CUSUM retardé de l'initié de l'exemple franchit-il son seuil~?
\end{enumerate}

\medskip\noindent\textbf{Partie IV --- Le verdict.}
\begin{enumerate}[resume]
\item Énoncer le \emph{résultat nommé}~: le nombre de jours avant qu'une exfiltration lente d'un initié placé soit signalée, à une fausse alerte par mois pour cent utilisateurs,
et le nombre de combinaisons de rôles toxiques dans la politique de la société avant et après la correction.
\item Qu'ajouteriez-vous d'abord aux détecteurs~?
\item Comment la zone de trading doit-elle être segmentée~?
\item Quelle est la leçon du hook pour toute société qui signe des transactions~?
\item En une phrase~: de quoi traite la sécurité dans une société de trading~?
\end{enumerate}
\end{problem}

\section{Questions d'entretien}

\begin{interviewq}[$\star$ \iqroles{développeur}]\label{iq:pl:security:1}
Où une clé secrète d'API de bourse doit-elle résider, et comment une stratégie doit-elle l'obtenir~?
\end{interviewq}

\begin{interviewq}[$\star\star$ \iqroles{développeur}]\label{iq:pl:security:2}
Quelles permissions une clé d'API de bourse utilisée pour le trading doit-elle avoir~?
\end{interviewq}

\begin{interviewq}[$\star\star$ \iqroles{développeur}]\label{iq:pl:security:3}
Comment détecteriez-vous un employé qui copie lentement, sur plusieurs mois, le code d'une stratégie~?
\end{interviewq}

\begin{interviewq}[$\star\star$ \iqroles{développeur}]\label{iq:pl:security:4}
Qu'est-ce que le zero trust, et que change-t-il au sein d'un réseau de trading~?
\end{interviewq}

\begin{interviewq}[$\star\star\star$ \iqroles{développeur}]\label{iq:pl:security:5}
Concevez le contrôle d'accès, la \omterm{def:pl:security:secrets}{gestion des secrets} et la surveillance d'une société de trading de cent personnes disposant de clés de plateformes et d'une trésorerie en crypto-actifs.
\end{interviewq}
